\begin{examplebox}
	\textbf{\textcolor{CExample}{例 1（基础）}}

	\textbf{\textcolor{CExample}{题目：}} 若正四面体的棱长为 $2$，求其外接球的半径 $R$。

	\textbf{\textcolor{CExample}{解题步骤：}}
	\begin{description}[style=nextline, leftmargin=2.8em, labelsep=0.5em, itemsep=0.45em]
		\item[\textbf{第一步（代入公式）：}] 直接使用正四面体外接球半径公式：
		      \[
			      R = \frac{\sqrt{6}}{4}a.
		      \]
		      代入 $a=2$，得
		      \[
			      R = \frac{\sqrt{6}}{4}\times 2.
		      \]
		      \par\textbf{理由：} 题目明确是正四面体，满足公式适用条件。

		\item[\textbf{第二步（计算数值）：}] 完成乘法运算：
		      \[
			      R = \frac{\sqrt{6}}{4}\times 2 = \frac{\sqrt{6}}{2}.
		      \]
		      \par\textbf{理由：} 分数乘法，约简系数。
	\end{description}

	\textbf{\textcolor{CExample}{关键结论：}} \textbf{$R = \frac{\sqrt{6}}{2}$。}

	\medskip
	\textbf{\textcolor{CExample}{例 2（稍变形）}}

	\textbf{\textcolor{CExample}{题目：}} 已知正四面体的高 $h = \sqrt{6}$，求其外接球半径 $R$。

	\textbf{\textcolor{CExample}{解题步骤：}}
	\begin{description}[style=nextline, leftmargin=2.8em, labelsep=0.5em, itemsep=0.45em]
		\item[\textbf{第一步（使用棱高关系）：}] 正四面体外接球半径与高满足
		      \[
			      R=\frac{3}{4}h.
		      \]
		      \par\textbf{理由：} 正四面体中，外接球球心在高线上，且到顶点的距离为高的 $\frac{3}{4}$。

		\item[\textbf{第二步（代入计算）：}] 将 $h=\sqrt{6}$ 代入：
		      \[
			      R=\frac{3}{4}\sqrt{6}=\frac{3\sqrt{6}}{4}.
		      \]
		      \par\textbf{理由：} 直接利用 $R=\frac{3}{4}h$ 比先求棱长更简洁。
	\end{description}

	\textbf{\textcolor{CExample}{关键结论：}} \textbf{$R = \frac{3\sqrt{6}}{4}$。}

	\medskip
	\textbf{\textcolor{CExample}{例 3（综合应用）}}

	\textbf{\textcolor{CExample}{题目：}} 若正四面体的外接球表面积为 $12\pi$，求该正四面体的棱长 $a$。

	\textbf{\textcolor{CExample}{解题步骤：}}
	\begin{description}[style=nextline, leftmargin=2.8em, labelsep=0.5em, itemsep=0.45em]
		\item[\textbf{第一步（求外接球半径）：}] 由球的表面积公式 $S = 4\pi R^2$ 得：
		      \[
			      12\pi = 4\pi R^2
			      \quad\Rightarrow\quad
			      R^2 = 3
			      \quad\Rightarrow\quad
			      R = \sqrt{3}.
		      \]
		      \par\textbf{理由：} 等式两边同除以 $4\pi$，再取算术平方根。半径为非负数，所以 $R=\sqrt{3}$。

		\item[\textbf{第二步（反求棱长）：}] 由
		      \[
			      R = \frac{\sqrt{6}}{4}a
		      \]
		      解出
		      \[
			      a = \frac{4R}{\sqrt{6}}.
		      \]
		      代入 $R=\sqrt{3}$，得
		      \[
			      a = \frac{4\sqrt{3}}{\sqrt{6}}
			      = \frac{4}{\sqrt{2}}
			      = 2\sqrt{2}.
		      \]
		      \par\textbf{理由：} 公式变形，化简时 $\frac{\sqrt{3}}{\sqrt{6}} = \frac{1}{\sqrt{2}}$。
	\end{description}

	\textbf{\textcolor{CExample}{关键结论：}} \textbf{棱长 $a = 2\sqrt{2}$。}
\end{examplebox}